\section{Transiciones}
\label{sec:transiciones}

Las transiciones controlan cómo se realiza el cambio de una fuente o composite a otro
en el flujo de programa. Voctomix 2.0 ofrece tres modos de transición que el realizador
activa desde la barra de herramientas de la interfaz gráfica.

\subsection{Cut (corte instantáneo)}

El botón Cut realiza una transición instantánea: en el siguiente fotograma
tras la pulsación, la fuente de programa pasa a ser la fuente de preview. No existe
ningún efecto intermedio. Es la transición estándar en producción broadcast: el
corte imita el montaje cinematográfico y resulta imperceptible cuando el
\textit{timing} es correcto.

\subsection{Trans (transición suave)}

El botón Trans realiza una transición animada entre el estado actual y el
nuevo estado de composite o fuente, interpolando los parámetros geométricos y de
transparencia durante un tiempo configurable (por defecto 750~ms, definido en la
sección \texttt{[transitions]} del fichero de configuración). La transición se implementa
modificando progresivamente los valores de los pads del compositor de GStreamer.

Voctomix 2.0 define múltiples trayectorias de transición según el estado de origen y
destino (por ejemplo, de \texttt{fs} a \texttt{pip} la transición pasa por el composite
intermedio \texttt{fs-pip} antes de llegar a \texttt{pip}), garantizando que la
animación sea visualmente coherente independientemente del cambio solicitado.

\subsection{Retake}

El botón Retake permite al realizador deshacer el último cambio de fuente o
composite y volver al estado inmediatamente anterior sin interrumpir la señal de
programa. Su funcionamiento es equivalente a un Cut inverso: en el siguiente fotograma,
el sistema intercambia programa y preview, restaurando la fuente que estaba en antena
antes del cambio. Esta función resulta especialmente útil cuando el realizador realiza
un corte erróneo y necesita corregirlo de inmediato sin generar un segundo corte
visible en la emisión.

\begin{figure}[H]
  \centering
  \fbox{\begin{minipage}{0.7\textwidth}\centering\vspace{1.2cm}
    \textit{[IMAGEN PENDIENTE: captura de los botones CUT, TRANS y RETAKE}\\
    \textit{en la barra de herramientas de voctogui.}\\[4pt]
    \textit{Guardar como \texttt{figuras/panel\_transiciones.png} y enviar.}
    \vspace{1.2cm}
  \end{minipage}}
  \caption{Botones de transición en la interfaz de Voctomix~2.0: CUT, TRANS y RETAKE.}
  \label{fig:panel_transiciones}
\end{figure}
